%==============
\begin{center}
{\LARGE \bf Non-Standard Applications
}
\end{center}
%==============
\bigskip

{\bf Organizer}
Eugenio Roanes-Lozano (eroanes@eucmos.sim.ucm.es)

\medskip

\bigskip

{\bf Description}
\medskip

Computer Algebra has traditionally been used in fields like Algebraic Geometry, Astronomy,
High Energy Physics, Automatic Theorem Proving in Geometry,... Nevertheless there are some
other fields such as Industrial Applications, Chemistry, Artificial Intelligence (for instance the
AISMC conference is entirely devoted to this topic),... where CA is been used now as a very useful
tool. The session tries to give an overview of some of these non-standard applications. Emphasis is
given in the originality. There will be a final colloquium about the topic. 

\bigskip
{\bf Talks}
\medskip

Date: July 18th (Thursday) 

\begin{description}
\item[14:00-14:25] \ \\
  Jacques Calmet\\
  {\bf Computer Algebra and Artificial Intelligence.} 
\item[14:25-14:50 ] \ \\
  Eugenio Roanes-Lozano, Luis M. Laita\\
  {\bf A Topology-Independent Model for Railway Interlocking Systems.} 
\item[14:50-15:15] \ \\
  Antonio Montes \\
  {\bf Algebraic Solution of the Load-Flow Problem for a four Nodes Electrical Network.
} 
\item[15:15-15:40] \ \\
  Antonio Montes\\
  {\bf Algebraic Solution of the Load-Flow Problem for a four Nodes Electrical Network.} 
\item[15:15-15:40] \ \\
  R.L. Akers, P. Baffes, E. Kant, C. Randall, S. Steinberg, R.L Young
\\
  {\bf A Problem Solving Environment for Numerical Partial
    Differential Equations.} 
\item[15:40-16:00] \ \\
  \\
  {\bf Coffee Break} 
\item[16:00-16:25] \ \\
  Shigekazu Nakagawa, Naoto Niki, Hiroki Hashiguchi
\\
  {\bf Computer Algebra application to the distribution of sample
    correlation  coefficient.} 
\item[16:25-16:50 ] \ \\
  Alejandro A.R. Trejo, Guillermo Fernandez Anaya
\\
  {\bf Regular Expressions Simplification.
} 
\item[16:50-17:15 ] \ \\
  Ross Taylor\\
  {\bf Thermodyanmics with Maple.} 
\end{description}


